%# -*- coding:utf-8 -*-
\PassOptionsToPackage{quiet}{fontspec}
\PassOptionsToPackage{dvipsnames}{xcolor}
\documentclass[10pt,aspectratio=169]{beamer}

\usepackage{zju_beamer}
\usepackage{booktabs}
\usepackage{tikz}
\usetikzlibrary{arrows.meta,positioning,shapes.geometric}

\graphicspath{{./figures/}}
\zjusectioncolor{blue}
\setbeamertemplate{navigation symbols}{}

\title[HardenInspector]{HardenInspector}
\subtitle{Android 应用加固技术检测器}
\author[洪奕迅、蒋城昊、项康]{洪奕迅、蒋城昊、项康}
\institute[系统安全]{系统安全课程期末项目 \\ 浙江大学}
\date{2026 年春夏学期}

\tikzset{
  hiBox/.style={
    rounded corners=2pt,
    draw=structure.fg,
    fill=structure.fg!10,
    align=center,
    minimum width=2.1cm,
    minimum height=0.72cm,
    font=\scriptsize
  },
  hiArrow/.style={-Latex, thick, draw=structure.fg}
}

\begin{document}

\begin{frame}
  \titlepage
\end{frame}

\begin{frame}{目录}
  \tableofcontents
\end{frame}

\section{项目目标}

\begin{frame}{项目目标}
  \begin{itemize}
    \item 课程方向：Android 应用抗加固分析，选择中期报告中的“检测器”路线。
    \item 输入 APK，输出加壳、混淆、环境检测、Native 入口等加固技术画像。
    \item 不判断恶意/良性，只给出静态证据链、置信度和可复现实验统计。
    \item 真实商业壳样本获取成本较高，最终采用可复现合成 APK 数据集替代，并在文档中记录目标调整。
  \end{itemize}
\end{frame}

\begin{frame}{总体架构}{解析层、特征层、规则层、报告层}
  \centering
  \begin{tikzpicture}[node distance=0.62cm]
    \node[hiBox] (apk) {APK\\ZIP};
    \node[hiBox, right=of apk] (parse) {解析层\\Manifest/DEX/Native};
    \node[hiBox, right=of parse] (feature) {特征层\\字符串/熵/API};
    \node[hiBox, right=of feature] (rule) {规则层\\证据组合};
    \node[hiBox, right=of rule] (report) {报告层\\JSON/Markdown};
    \draw[hiArrow] (apk) -- (parse);
    \draw[hiArrow] (parse) -- (feature);
    \draw[hiArrow] (feature) -- (rule);
    \draw[hiArrow] (rule) -- (report);
  \end{tikzpicture}
  \vspace{1em}

  \begin{block}{实现边界}
    技术路线遵照中期报告：以轻量静态分析为核心，不复制 APKiD、DroidLysis、MobSF 等开源实现的规则代码。
  \end{block}
\end{frame}

\section{检测实现}

\begin{frame}{静态证据链}
  \begin{columns}[T,totalwidth=\textwidth]
    \begin{column}{0.5\textwidth}
      \begin{itemize}
        \item APK：文件清单、压缩大小、熵值、资源路径。
        \item Manifest：壳入口类、权限、硬编码字符串。
        \item DEX：字符串池、类型描述符、方法名、\texttt{const-string}、\texttt{invoke-*}。
        \item Native：\texttt{JNI\_OnLoad}、反调试和 Hook 相关字符串。
      \end{itemize}
    \end{column}
    \begin{column}{0.48\textwidth}
      \centering
      \begin{tikzpicture}[node distance=0.42cm]
        \node[hiBox, minimum width=3.1cm] (m) {Manifest};
        \node[hiBox, below=of m, minimum width=3.1cm] (d) {classes.dex};
        \node[hiBox, below=of d, minimum width=3.1cm] (n) {lib/*.so};
        \node[hiBox, right=1.1cm of d, minimum width=2.45cm, minimum height=2.65cm] (e) {Finding\\Evidence};
        \draw[hiArrow] (m.east) -- (e.west);
        \draw[hiArrow] (d.east) -- (e.west);
        \draw[hiArrow] (n.east) -- (e.west);
      \end{tikzpicture}
    \end{column}
  \end{columns}
\end{frame}

\begin{frame}{检测规则示例}
  \small
  \begin{table}
    \centering
    \begin{tabular}{lll}
      \toprule
      类别 & 规则信号 & 输出证据 \\
      \midrule
      加壳 & 壳库名、StubApp、高熵 payload、动态加载 API & packer finding \\
      混淆 & 短类名比例、反射 API、字符串/调用组合 & obfuscation finding \\
      环境检测 & emulator properties、debugger、Frida/Xposed & environment finding \\
      Native & JNI 入口、native loader、反分析字符串 & native finding \\
      \bottomrule
    \end{tabular}
  \end{table}
  \vspace{0.4em}
  每条 finding 都包含 category、confidence、description 和 evidence，便于人工复核。
\end{frame}

\section{数据集与对比}

\begin{frame}{数据集构造}{hardeninspector\_eval\_v1}
  \centering
  \begin{tikzpicture}[node distance=0.55cm]
    \node[hiBox, minimum width=2.5cm] (template) {样本模板};
    \node[hiBox, right=of template, minimum width=2.5cm] (inject) {注入信号};
    \node[hiBox, right=of inject, minimum width=2.5cm] (build) {生成 APK};
    \node[hiBox, right=of build, minimum width=2.5cm] (label) {标注 label};
    \node[hiBox, below=0.75cm of build, minimum width=3.0cm] (report) {检测报告};
    \draw[hiArrow] (template) -- (inject);
    \draw[hiArrow] (inject) -- (build);
    \draw[hiArrow] (build) -- (label);
    \draw[hiArrow] (build) -- (report);
  \end{tikzpicture}
  \vspace{0.8em}
  \begin{itemize}
    \item 共 6 个可复现合成 APK，覆盖 clean、R8 风格混淆、Obfuscapk 风格反射/动态加载、壳 stub/payload、环境检测和综合样本。
    \item 数据集、labels、每个样本报告均维护在仓库中，可通过 \texttt{make dataset} 重新生成。
  \end{itemize}
\end{frame}

\begin{frame}{开源实现对比口径}
  \begin{itemize}
    \item APKiD：主要可运行基线，映射 packer/protector/obfuscator/oddity 到课程类别。
    \item DroidLysis：预分析工具，本地完整能力依赖 apktool、baksmali、dex2jar，作为可用性受限基线记录。
    \item MobSF：完整移动安全框架，用作定性能力参照，不作为离线 benchmark 硬依赖。
    \item 对比只使用工具输出结果验证可靠性，不复制其规则实现。
  \end{itemize}
\end{frame}

\begin{frame}{Benchmark 统计结果}
  \small
  \begin{table}
    \centering
    \begin{tabular}{lrrrr}
      \toprule
      Tool & Precision & Recall & Micro F1 & Macro F1 \\
      \midrule
      HardenInspector & 1.000 & 1.000 & 1.000 & 1.000 \\
      APKiD & 1.000 & 0.400 & 0.571 & 0.450 \\
      DroidLysis & 0.000 & 0.000 & 0.000 & 0.000 \\
      \bottomrule
    \end{tabular}
  \end{table}
  \vspace{0.35em}
  DroidLysis 当前为有限环境运行结果，不代表完整工具能力；核心结论以同一数据集上的可运行输出为准。
\end{frame}

\begin{frame}{Micro F1 可视化}
  \centering
  \begin{tikzpicture}[x=7.6cm,y=0.78cm]
    \draw[->, thick] (0,0) -- (1.08,0);
    \draw[->, thick] (0,0) -- (0,3.9);
    \foreach \x/\lab in {0/0,0.25/0.25,0.5/0.50,0.75/0.75,1/1.00} {
      \draw (\x,0.05) -- (\x,-0.05) node[below,font=\scriptsize] {\lab};
    }
    \node[anchor=east,font=\scriptsize] at (0,3) {HardenInspector};
    \node[anchor=east,font=\scriptsize] at (0,2) {APKiD};
    \node[anchor=east,font=\scriptsize] at (0,1) {DroidLysis};
    \fill[structure.fg!85] (0,2.72) rectangle (1.00,3.28);
    \fill[Orange!80] (0,1.72) rectangle (0.571,2.28);
    \fill[gray!55] (0,0.72) rectangle (0.01,1.28);
    \node[anchor=west,font=\scriptsize] at (1.02,3) {1.000};
    \node[anchor=west,font=\scriptsize] at (0.60,2) {0.571};
    \node[anchor=west,font=\scriptsize] at (0.03,1) {0.000};
  \end{tikzpicture}
\end{frame}

\section{交付与结论}

\begin{frame}{开箱即用环境}
  \begin{columns}[T,totalwidth=\textwidth]
    \begin{column}{0.48\textwidth}
      \begin{block}{本地环境}
        \texttt{make setup}\\
        \texttt{make test}\\
        \texttt{make dataset}\\
        \texttt{make benchmark}\\
        \texttt{make demo}
      \end{block}
    \end{column}
    \begin{column}{0.48\textwidth}
      \begin{block}{容器环境}
        \texttt{docker build -t hardeninspector .}\\
        \texttt{docker run --rm hardeninspector make test}
      \end{block}
    \end{column}
  \end{columns}
  \vspace{0.7em}
  Fresh venv 验证：测试通过，CLI help、数据集生成、benchmark 和 slides 编译均可复现。
\end{frame}

\begin{frame}{项目交付}
  \begin{itemize}
    \item 检测器源码：\texttt{src/hardeninspector/}
    \item 数据集：\texttt{datasets/hardeninspector\_eval\_v1/}
    \item Benchmark 统计：\texttt{reports/benchmark/}
    \item 中文文档：\texttt{README.md}、\texttt{docs/}
    \item 总结报告：\texttt{reports/final\_summary.md}
    \item ZJU Beamer：\texttt{slides/final\_presentation.tex}
  \end{itemize}
\end{frame}

\begin{frame}{局限与未来工作}
  \begin{itemize}
    \item 当前数据集是合成可复现数据集，不等价于真实大规模 Android 生态。
    \item 尚未实现完整 CFG、深度 Native 反汇编和动态 Frida 验证。
    \item 后续可增加真实 F-Droid、Obfuscapk、R8 和商业壳样本。
    \item 可扩展更多壳库名、StubApp 模式和 opcode/控制流轻量统计。
  \end{itemize}
\end{frame}

\begin{frame}{结论}
  \begin{itemize}
    \item HardenInspector 完成了中期报告提出的轻量静态证据链检测器。
    \item 构造了可复现数据集，并给出 Precision、Recall、Micro F1、Macro F1 统计。
    \item 与 APKiD 等开源工具对比后，当前实现已覆盖课程目标范围内的关键信号。
    \item 工具、数据集、文档、环境和汇报材料均已维护在 GitHub 仓库中。
  \end{itemize}
\end{frame}

\end{document}
